\documentclass[12pt, a4paper]{article}

\usepackage[margin=1.2in]{geometry}
\usepackage{setspace}
\usepackage{microtype}
\usepackage{hyperref}
\usepackage{amsmath}
\usepackage{amssymb}
\usepackage{csquotes}
\usepackage{fontenc}
\usepackage{inputenc}
\usepackage{parskip}
\usepackage{titlesec}
\usepackage{fancyhdr}
\usepackage{abstract}

\onehalfspacing

\hypersetup{
    colorlinks=true,
    linkcolor=black,
    urlcolor=black,
    citecolor=black
}

\titleformat{\section}
  {\normalfont\large\bfseries}{\thesection.}{0.75em}{}
\titlespacing*{\section}{0pt}{2ex plus 0.5ex minus 0.2ex}{1ex plus 0.2ex}

\pagestyle{fancy}
\fancyhf{}
\fancyhead[R]{\thepage}
\fancyhead[L]{\textit{The Hollow Network}}
\renewcommand{\headrulewidth}{0.4pt}

\title{\textbf{The Hollow Network}\\[0.5em]
\large Synthetic Sociality, Metric Misalignment, and the Politics of Legibility in Algorithmically Mediated Environments}
\author{Flyxion}
\date{}

\begin{document}

\maketitle
\thispagestyle{empty}

\begin{abstract}
\noindent This essay traces the emergence of what it calls \emph{synthetic sociality}---communicative environments in which the proportion of automated, fabricated, or algorithmically amplified activity reaches a threshold at which ordinary users can no longer reliably distinguish human interaction from simulation. Beginning with the phenomenological experience of impersonation, bot networks, and manufactured content on large social platforms, the essay progressively widens its explanatory frame. It argues that synthetic sociality is not primarily the product of malicious actors operating against the interests of platforms, but an emergent property of metric-optimization systems that treat automated engagement as functionally equivalent to genuine participation. It further argues that describing platforms as unified rational maximizers misrepresents the organizational dynamics through which technological environments are actually produced. Drawing on historical comparison, the essay considers whether present conditions represent a novel failure mode or merely the latest iteration of longstanding anxieties about mediated communication. It concludes that the central problem is not the existence of synthetic activity but its opacity and involuntary character, and that the appropriate design response centers on informational legibility, provenance, and the transformation of synthetic interaction from a hidden structural condition into a transparent and consensual one.
\end{abstract}

\newpage
\setcounter{page}{1}

\section{The Perceptual Contradiction}

There is a peculiar dissonance at the center of contemporary social media. The largest platforms are also, by their own account, among the most sophisticated artificial intelligence laboratories in existence. Their research divisions publish widely in machine learning, their infrastructure processes billions of interactions daily, and their public communications consistently frame artificial intelligence as a transformative project: a technology that will improve content moderation, personalize discovery, and deepen the quality of human connection across networks of unprecedented scale.

Yet the everyday experience of using these platforms increasingly consists of something quite different. An ordinary user logging into a major social network in 2024 is likely to encounter impersonation accounts cloning the profile photographs and names of people they know, coordinated bot networks producing and amplifying content at velocities no human community could sustain, and entire pages or channels composed of synthetic or recycled video assembled by automated pipelines optimized not for meaning but for engagement signals. The friend request from a duplicate account, the adoption listing that collects a deposit and then vanishes, the comment thread where suspiciously similar accounts converge on an identical talking point---these are not marginal experiences but routine features of the social environment these platforms have built.

This creates an initial puzzle. If a company invests billions in frontier AI research, why does its platform feel crowded with crude identity scams? If recommendation algorithms operate at a level of sophistication that can model individual attention patterns across hundreds of millions of users, why cannot the same organization reliably detect when an account has cloned a photograph from another profile and is soliciting money from that person's contacts?

The easy answer---that these are simply different technical problems, and that malicious actors are an inevitable presence in any large open system---is not wrong, but it is insufficient. It treats the distribution of engineering effort as though it were a natural phenomenon, a consequence of difficulty rather than of institutional choice. The harder question is why certain capabilities receive resources and others do not, and what those priorities reveal about the actual objectives these organizations are designed to pursue.

\section{The Infrastructure Alibi}

The initial explanation most commonly offered for the proliferation of synthetic activity on social platforms is structural. Large-scale automation infrastructure, the argument goes, is inherently dual-use. The same pipelines that allow legitimate applications---content recommendation, advertising targeting, automated customer service---necessarily lower the barrier for actors pursuing illegitimate ends. A platform that makes it easy to distribute video at scale makes it equally easy for a reel farm to flood the recommendation system with manufactured content. A network that allows rapid account creation enables both authentic participation and mass-produced impersonation.

This explanation is accurate as far as it goes. The technical infrastructure of large platforms does not distinguish, at a fundamental level, between human activity and automated mimicry of human activity. Both produce signals---clicks, watch time, shares, follow actions---that the system is built to process and respond to. There is no architectural feature of a recommendation pipeline that inherently privileges genuine human intention over automated approximation of it.

But the structural explanation carries an implicit exculpatory logic. By framing synthetic activity as an unfortunate side effect of general-purpose infrastructure, it positions the platform as a neutral party whose tools have been turned against it by external malicious actors. The platform, in this framing, is struggling with a problem that has been imposed on it from outside---adversaries exploiting infrastructure that was designed for better purposes.

This framing has rhetorical utility. It allows companies to speak simultaneously about their commitment to safety and their investment in AI without confronting the possibility that these two commitments might be in tension. It also determines which interventions appear natural: if the problem is external adversaries exploiting legitimate infrastructure, the solution is better detection and enforcement---a security problem, technical in character, requiring engineering resources but not organizational reflection.

What the framing does not address is the prior question of why the infrastructure was built without authenticity requirements in the first place, and why it has remained without them as the scale and sophistication of synthetic activity has become increasingly legible.

\section{Metric Indifference and Its Beneficiaries}

A stronger interpretation of the phenomenon requires attending more carefully to what social platforms actually optimize and for whom. The dominant metric across large consumer platforms is some variant of engagement: click-through rates, watch time, shares, daily active users, time on platform. These indicators began as proxies for something that seemed genuinely meaningful---evidence that people were finding value in the network, that content was resonating, that the platform was doing its job of connecting people with things they wanted to see.

The moment these proxies became the primary optimization target of both internal systems and external reporting, however, they acquired a different character. Goodhart's Law---the principle that when a measure becomes a target, it ceases to be a reliable measure of what it originally represented---describes the resulting dynamic with precision. Once engagement metrics entered the objective functions of recommendation algorithms, quarterly reports, and engineering performance reviews, the system's behavior reoriented around the production of those metrics rather than around whatever underlying quality they were originally meant to indicate.

We can formalize the dynamic briefly. Let $E(A)$ denote the engagement produced by activity $A$ in the platform environment. The optimization problem the system actually solves is:

\[
\max_{A} \; E(A) \quad \text{subject to operational constraints}
\]

The critical observation is that if authenticity---the property that $A$ represents genuine human expression or communication---does not appear in $E$, then it is invisible to the objective function. An automated account that watches a video long enough to register a completion event, triggers a recommendation to another account, and generates a share action produces exactly the same signal as a human performing identical interactions. From the perspective of the optimization system, these are equivalent.

This equivalence is not merely a technical artifact; it has economic consequences. Advertising markets are primarily concerned with impressions and interaction signals. An automated impression that generates a measurable interaction is, from the standpoint of the advertising transaction, indistinguishable from a human one---provided the fraud does not reach a scale that advertisers can detect and penalize. At moderate levels, automated activity can actually stabilize platform metrics, filling engagement gaps when human participation fluctuates and providing consistent throughput for recommendation pipelines.

The implication is uncomfortable but analytically unavoidable. Platforms do not merely tolerate synthetic activity because it is technically difficult to eliminate. They operate in conditions where synthetic activity produces measurable signals that their business model treats as positive---at least below whatever threshold would attract regulatory scrutiny or advertiser defection. The relationship between the platform and the synthetic activity is therefore not cleanly adversarial. It is, at minimum, one of mutual indifference, and in some respects one of mutual benefit.

This does not require deliberate deception. It can arise from something more mundane: a system designed to maximize measurable signals, encountering a class of activity that produces those signals, and failing to develop the capacity to distinguish it from the activity the signals were meant to represent. The absence of malice does not, however, diminish the structural condition. An optimization system that is indifferent to authenticity will produce environments in which authenticity is not protected, regardless of what the organization professes in its public communications.

\section{Against the Unified Maximizer}

There is a danger in the analysis developed so far, which is that it generates too clean a picture. Writing the system as $\max E(A)$ implies the existence of a coherent agent with a well-defined objective function. Large technology companies are not like this. They are bureaucratic organizations of tens or hundreds of thousands of people, distributed across divisions that have different mandates, different time horizons, and sometimes directly competing incentives.

An engineering team building recommendation infrastructure is rewarded for model performance on engagement metrics. A trust-and-safety team is rewarded for the volume of policy violations detected and removed. A product team is rewarded for user growth and retention numbers. A policy team responds to regulatory pressure and manages relationships with governments and civil society organizations. Each of these teams optimizes for a locally defined objective. The interaction of these locally optimized systems produces the overall behavior of the organization---but that behavior does not have the coherence one would expect if a single agent were solving a single problem.

This matters because it changes the diagnosis of why synthetic activity persists. The answer may not be that the company has made a strategic decision to tolerate bots because they inflate metrics. It may be something more dispersed: that the team with the mandate to address coordinated inauthentic behavior has insufficient resources relative to the teams building the infrastructure that enables it; that detection of sophisticated synthetic networks requires capabilities that cross organizational boundaries and therefore fall into gaps between divisions; that the metrics by which senior leadership evaluates platform health do not include indicators of synthetic activity that would make its prevalence visible; or simply that the problem is technically difficult and the organization's attention is directed elsewhere by immediate competitive pressures.

These organizational dynamics are not exculpatory---they explain rather than excuse---but they do have implications for how change might occur. An organization that has made a strategic choice to prioritize engagement over authenticity might be moved by changes in regulation, advertising market pressure, or public discourse to shift that choice. An organization whose behavior toward synthetic activity emerges from distributed organizational dysfunction requires different interventions: changes to internal measurement systems, resource allocation, and promotion criteria that bring the stated commitment to authenticity into alignment with the incentive structures that actually govern day-to-day decisions.

The institutional explanation also draws attention to a pattern that appears across many organizations managing complex infrastructures at scale: the gap between what an organization publicly commits to and what its internal incentive architecture actually rewards tends to be visible precisely in the areas where commitment would be costly. Eliminating synthetic activity aggressively might shrink reported user bases, reduce engagement metrics, and create friction in content distribution pipelines---all of which would appear as negative signals in the reporting structures by which large platforms are evaluated. This is not a conspiracy; it is the ordinary operation of institutional incentives, and it is why commitments to authenticity that are not embedded in measurable internal targets tend to remain symbolic.

\section{The Externalization of Cost}

The analysis of platform incentives acquires a sharper moral character when considered from the perspective of the people who bear the consequences of synthetic activity. The asymmetry between those who create fraud and those who experience it is not accidental. It is structured by the same institutional indifference that allows synthetic activity to persist within the optimization environment.

Consider the range of fraud enabled by the conditions described above. Impersonation accounts clone a person's photograph, name, and fragments of their public posts to create convincing duplicates, which are then used to solicit money from the original person's contacts or to distribute fraudulent links to their audience. Fake rehoming pages for animals accept deposits from people seeking to adopt pets and then disappear, with the operator moving to a new page or a different regional community group. Rental scams collect deposits on apartments that do not exist or are not available for rent, exploiting the trust that platform community membership implies. Investment fraud accounts build apparent legitimacy through stolen identity signals before promoting schemes to extract money from followers.

These are not sophisticated attacks. They exploit elementary features of platform design: the ease of account creation without identity verification, the absence of meaningful history or provenance signals for new accounts, the capacity to operate across jurisdictions and communities without accountability to any of them, and the availability of algorithmic amplification that places new content in front of large audiences before any reputation has been established.

The cost of this fraud is borne almost entirely by individuals. A person who has lost a deposit to a rental scam must determine, individually, whether and how to pursue recovery---a process that is time-consuming, often fruitless, and emotionally taxing. A person whose identity has been impersonated must navigate reporting systems that may respond slowly or inconsistently, while their contacts continue to receive fraudulent messages from the duplicate account. The platform removes an account if and when it violates policy; the network behind it persists and generates new accounts.

This pattern reflects a broader principle about the distribution of costs in systems with weak accountability mechanisms. In offline contexts, fraud encounters friction: social networks impose reputational costs on deception that is discovered, legal systems create deterrence, and the geographic and social proximity of fraudulent activity to its victims creates pressure toward accountability. None of these friction mechanisms operate effectively in large digital environments with low barriers to account creation and identity change.

The jarring quality of this situation---and it is genuinely jarring---is the contrast between the organizational resources directed toward capability development and the relative thinness of basic protective infrastructure. It is not primarily a technical problem that a platform investing at scale in AI systems cannot build reliable identity verification, enforce against repeat fraudulent operators, or surface provenance information for accounts and content. It is an institutional problem: the investments that would address these failures do not register as improvements in the metrics by which these organizations are governed and evaluated.

One might observe that analogous dynamics existed long before digital platforms. Gossip and rumor in small communities could damage reputations without recourse; fraudulent operators in physical marketplaces found ways to exploit trust and disappear. The platforms are not inventors of social deception. What they have done is remove the friction that historically constrained its scale. In a small community, the same social network that transmitted rumor also limited its reach and eventually created accountability. On a platform serving hundreds of millions of users across thousands of communities, an operator can exploit each community in sequence, moving faster than any single community can respond, with no mechanism for coordination among victims that would produce collective pressure on the institution hosting the activity.

\section{The Historical Mirror}

Before concluding that contemporary platforms represent a wholly unprecedented failure, it is worth submitting the diagnosis to historical comparison. Anxieties about the authenticity of mediated communication have accompanied nearly every significant communication technology. The introduction of print raised concerns that the proliferation of texts would overwhelm the reader's capacity for critical evaluation. Broadcasting was accused of manufacturing false consensus, homogenizing culture, and replacing genuine political participation with the appearance of it. Early internet forums generated extensive discussion about the unreliability of anonymous communication and the impossibility of verifying the identity or sincerity of interlocutors.

In each case, critics worried that technological mediation would displace authentic human contact with a cheaper substitute---that the medium would produce environments populated by representations and performances rather than genuine persons. And in each case, the prediction was partially vindicated and partially refuted: new media did introduce new forms of manipulation and inauthenticity, and human communities found ways to develop new norms, literacies, and institutions that allowed meaningful interaction to persist alongside the degraded forms.

The question for contemporary platforms is whether they represent merely the latest iteration of this recurrent anxiety or whether they introduce failure modes that are qualitatively different from what preceded them. There are grounds for arguing that some features of the current situation are genuinely novel. Earlier mass media manipulated at scale but lacked personalization: a propaganda campaign addressed the same message to a mass audience. Contemporary platforms combine algorithmic personalization, granular data about social relationships, and the capacity to generate synthetic identities that can operate inside individual social networks rather than broadcasting to an undifferentiated public.

An impersonation account that clones a specific person's identity and operates within that person's actual social graph---sending messages to their actual friends and colleagues---is categorically different from a mass-media propaganda campaign. It exploits the intimacy of personal relationship rather than the diffuse attention of a public audience. It is harder to detect because it does not feel like mass communication; it feels like a message from someone the recipient knows. And it can be generated at scale with minimal cost, targeting thousands of individual social contexts simultaneously in ways that no earlier technology permitted.

This qualitative distinction matters because it determines what forms of literacy and institutional response are appropriate. Teaching people to be skeptical of anonymous strangers on early internet forums was a reasonable response to that environment. Teaching people to suspect messages from accounts that appear to be their friends and family members, while not technically unreasonable as advice, points toward the collapse of a social environment rather than its maintenance.

\section{The Complexity of Desire}

Any account of what has gone wrong with social platforms must grapple honestly with the complexity of what users actually want. The analysis so far has implicitly treated the preference for authentic human interaction as given, and the proliferation of synthetic activity as a deviation from that preference that users would eliminate if given the choice. This is not obviously correct.

Many people engage willingly and knowingly with automated or semi-automated communicative actors. Parasocial relationships with content creators who interact with audiences through algorithmically optimized posting schedules, brand accounts that communicate through automated responses, and entertainment streams produced by coordinated content teams are all widely chosen forms of digital interaction. They are often attractive precisely because they require no reciprocity: the engagement is one-directional, the social obligation minimal, and the experience of connection available without the cost of genuine relationship.

This complexity does not undermine the critique of involuntary synthetic activity, but it refines it. The problem users experience is not simply the existence of automated actors in their social environments. It is the specific character of that automation: deceptive, opaque, and non-consensual. The impersonation account is objectionable not because it is automated but because it claims to be something it is not, exploiting the trust that authentic identity implies. The coordinated bot network is problematic not because it involves multiple accounts but because it simulates the appearance of independent consensus that does not exist. The rental scam is harmful not because it uses digital tools but because it exploits the informational asymmetry those tools create.

The distinction between consensual and non-consensual synthetic interaction is therefore crucial. A user who knowingly follows an automated entertainment account is making a choice about how to allocate their attention. A user who receives a message from what appears to be their colleague and is in fact an impersonation account has had that choice made for them---and made against their interests.

\section{The Opacity Problem}

This reframing---from the problem of inauthenticity to the problem of opacity---allows a more precise diagnosis of what has gone wrong. The question is not whether social platforms should contain automated actors, synthetic content, or algorithmically mediated interaction. In some sense they cannot avoid containing all of these. The question is whether the nature of the actors and the status of the content are legible to users.

Opacity operates at several levels. At the level of identity, accounts that represent automated systems or organized promotional campaigns are often indistinguishable from accounts representing individual humans, because the platform provides insufficient provenance information to make the distinction. At the level of content, media that has been synthetically generated or repurposed from other contexts frequently circulates without indicators of origin. At the level of behavior, coordinated networks of accounts that simulate independent opinion formation are designed to be invisible as networks---their coordination is the feature that makes them effective, and the platform's infrastructure does not surface that coordination to users encountering their output.

The consequences of opacity for trust are significant. When users cannot determine whether the account sending them a message represents a person they know, an automated system, or an organized fraud network, the epistemic status of every interaction is degraded. The rational response is generalized suspicion---a posture that damages genuine communication alongside fraudulent simulation. The cost of this degradation falls unevenly: sophisticated users with the time and resources to investigate suspicious activity bear it differently than users who lack those capacities.

The burden of verification has been externalized to individual users in a system that provides them with insufficient tools to perform it. This is the institutional failure at the center of the problem. A communication infrastructure operating at the scale of the major social platforms is not simply a marketplace of ideas in which users are responsible for their own epistemic hygiene. It is an environment whose design choices determine what information is available to whom, and those choices currently distribute informational disadvantage downward---toward users who are least equipped to compensate for it.

\section{Legibility as Infrastructure}

The appropriate response to the opacity problem is not to attempt the elimination of synthetic activity from social environments---an impossible goal, and arguably not a desirable one given that many forms of automation serve genuine user interests. The appropriate response is to treat informational legibility as infrastructure: a basic feature of the communicative environment that the platform is responsible for providing.

Concretely, this means several things. Identity provenance should be surfaced rather than obscured: accounts representing individuals, organizations, automated systems, or coordinated campaigns should carry indicators that allow users to understand what kind of communicative actor they are engaging with. This does not require the elimination of pseudonymity, which serves legitimate purposes; it requires distinguishing between human-controlled pseudonymous accounts and non-human or coordinated actors. Content origin should be similarly legible: media that has been synthetically generated, repurposed from other contexts, or distributed by automated pipelines should carry provenance indicators that are visible to the user encountering it.

Enforcement against repeat fraudulent operators should treat the underlying network, not merely the individual account, as the unit of accountability. A scammer who creates a new page after having a previous one removed has not been meaningfully deterred; the infrastructure has simply absorbed the cost of their continuation. Effective enforcement requires coordinating across account histories, behavioral patterns, and network connections in ways that make the operational friction of fraud substantially higher than the current environment imposes.

User controls should allow people to choose the degree of synthetic interaction they wish to engage with. A user who prefers to see only content from verified human accounts should have mechanisms to accomplish this. A user who is comfortable engaging with automated entertainment accounts should be able to do so knowingly. The goal is not to homogenize the social environment but to transform the current condition---in which synthetic interaction is an ambient, hidden structural feature---into one in which it is a transparent option.

These interventions require institutional commitment that goes beyond stated values. They require that the measurement systems, promotion criteria, and resource allocation decisions inside large platforms treat legibility as a performance indicator in the same way that engagement is currently a performance indicator. They require that the costs of fraud---which are currently externalized to individual users---be recognized as costs attributable to platform design, and that accountability for those costs be located at the institutional level where the relevant design choices are made.

\section{Conclusion: The Redesign of Trust}

The experience of contemporary social media that motivates this essay---fake accounts, synthetic reels, adoption scams, impersonation fraud---is best understood not as the failure of technology to catch up with malicious actors but as the emergent property of an optimization system that is indifferent to authenticity and organized in ways that externalize the cost of that indifference onto ordinary users.

This diagnosis requires widening the frame beyond technical analysis. The problem is economic, in that advertising markets reward signals that automated activity can produce as readily as human activity. It is organizational, in that the institutional structures of large platforms distribute responsibility for synthetic activity across divisions with incompatible incentives, producing persistent gaps in which fraud operates without meaningful accountability. It is historical, in that the current condition extends longstanding dynamics of mediated communication while introducing genuinely novel failure modes at the intersection of personalization, social-graph data, and synthetic identity generation. And it is political, in that the costs of the current condition are distributed in ways that track existing asymmetries of power and resource---with ordinary users bearing the burden of a degraded informational environment that institutions have the capacity but not the incentive to repair.

The path forward does not lie in nostalgia for a fully authentic social environment that may never have existed. Human communities have always contained performance, strategic self-presentation, organized influence, and deliberate manipulation. What distinguishes the present situation is not the presence of these dynamics but the collapse of the friction that historically constrained them and the opacity that prevents users from calibrating their responses to them.

Rebuilding trust in digital communicative environments therefore requires neither the elimination of synthetic activity nor a retreat from networked sociality. It requires the construction of informational legibility as a public good: the design of systems in which the provenance, nature, and coordination of communicative signals are visible to those who encounter them. In such an environment, synthetic interaction would remain available as a choice rather than being imposed as a condition. The difference between those two things is the difference between a communicative environment that users can navigate with understanding and one that they must inhabit with ambient suspicion.

The challenge is institutional before it is technical. The capacity to build more legible communicative infrastructures exists. Whether the organizations that control those infrastructures can be moved---through regulation, market pressure, or internal reform---to treat legibility as a core commitment rather than a rhetorical one is the central political question that the analysis developed here cannot answer, but cannot avoid posing.

\end{document}
